% Generated by Listas/make_statements.py from pratica-lecture-04.md, item 3. Do not edit.
\textbf{Source.} Gibbons (1992), \emph{Game Theory for Applied Economists}, Ch. 2 Problems (Section 2.2), Problem 2.6, p.~133 {[}PDF 140{]}. Reproduced verbatim. Answer in the source: none (Gibbons prints no answers).



\begin{quote}

2.6. Three oligopolists operate in a market with inverse demand given by \(P(Q) = a - Q\), where \(Q = q_1 + q_2 + q_3\) and \(q_i\) is the quantity produced by firm \(i\). Each firm has a constant marginal cost of production, \(c\), and no fixed cost. The firms choose their quantities as follows: (1) firm 1 chooses \(q_1 \ge 0\); (2) firms 2 and 3 observe \(q_1\) and then simultaneously choose \(q_2\) and \(q_3\), respectively. What is the subgame-perfect outcome?

\end{quote}

